\section*{Reviewer-requested clarifications and resolution of contrary evidence}
\textbf{Scope of the claim.} The experiments do not establish that equality of protected-group means implies fairness or distributional indistinguishability. The strongest mathematical statement is narrower: for the training-fitted contrast $d=\mu_1-\mu_0$, Fisher direction $w=S_W^{-1}d$, rank-one operator $P=dw^T/(w^Td)$, and midpoint $\mu_f=(\mu_0+\mu_1)/2$, the unconditional transformation $T_1(z)=z-P(z-\mu_f)$ maps the two fitted means to the same midpoint. Since $\mu_0-\mu_f=-d/2$, $\mu_1-\mu_f=d/2$, and $Pd=d$, we obtain $T_1(\mu_0)=\mu_f$ and $T_1(\mu_1)=\mu_f$. The guarantee concerns the first moment only. The operator is generally an oblique rank-one projector because $P^2=P$ but $P^T\neq P$ unless the Fisher and mean-separation directions coincide.

\textbf{Contrary experiment 1: no generic interior optimum.} The independent mean-line trajectory was required to test whether an intermediate intervention strength could dominate full equalization. It did not. Protected AUC changed from 0.9447 at $\alpha=0$ to 0.7914 at $\alpha=.5$ and 0.5000 at $\alpha=1$, while task macro-F1 changed from 0.7024 to 0.7429 and 0.7500. Under the frozen decision rule, the strongest fairness endpoint was therefore the full intervention. The interior-optimum hypothesis is explicitly falsified and is not retained as a contribution.

\textbf{Contrary experiment 2: mean equality is not distributional equality.} At full mean-line strength, the test-set mean gap was still 0.4476 and symmetric diagonal-Gaussian KL was 0.3795. Thus a zero training first-moment contrast cannot be interpreted as equality of covariance or higher-order structure. This negative result directly limits the interpretation of LEACE-CL and prevents the paper from using ``indistinguishable'' as a synonym for ``equal means''.

\textbf{Contrary experiment 3: nonlinear information remains after canonical LEACE.} In the clean-room matrix, nonlinear protected AUC decreased from 0.7218 to 0.5573, rather than to 0.5. We therefore report canonical LEACE as reducing recoverability in the tested probes, not as eliminating all protected information. The residual is a falsification of the stronger complete-erasure interpretation under the nonlinear stress test.

\textbf{Contrary experiment 4: sparse intervention does not outperform full LEACE.} Full LEACE reaches protected AUC 0.5000 with task AUC 0.8607, while the representative 60\%/$\lambda=1$ sparse point reaches protected AUC 0.7790 with task AUC 0.8602. The sparse point is consequently not presented as superior. Its role is to expose an interpretable operating frontier when intervention coverage itself is considered a resource. If a point was chosen after inspecting held-out outcomes, it is descriptive rather than confirmatory.

\textbf{Protocol correction.} The mean-line experiment conditions on the protected group and therefore cannot be treated as a direct empirical deployment test of the unconditional LEACE-CL operator. It is retained as a mechanism-level first-moment diagnostic. LEACE-CL itself is evaluated under a separate unconditional, train-fitted protocol. The real-data MILK10k gate is not promoted until its workflow artifact is independently audited.

\textbf{Reproducibility and negative evidence.} Every promoted or exploratory result is linked to an experiment specification, workflow/run, configuration, artifact, metrics, and falsification status. A historical helper that did not implement canonical LEACE is excluded from LEACE claims. Synthetic results are labelled exploratory. The paper therefore treats failed hypotheses, residual predictability, and incomplete real-data validation as evidence that constrains the conclusion rather than as omissions.
